% ==================== 第四章：模型综合评价与对比分析 ====================
\section{模型综合评价与对比分析}
\label{sec:evaluation}

本章在统一的数据口径与实验流程下，对第一章（确定性）、第二章（鲁棒优化）与第三章（滚动 + CVaR）三类方案进行综合评价与对比。我们重点关注收益分布尾部、风险—收益权衡、资源利用与结构稳定性，以及价格—供给—容量反馈下的管理含义。

\subsection{实验设置}
为保证可比性，沿用全文统一的数据与规模设定：作物集 $\mathcal K=6$、地块 $\mathcal I=24$、季节 $\mathcal S=3$、评估周期 $T=7$ 年。单产按对数正态波动，价格采用线性或对数线性的价格—供给反馈，并显式考虑销量上限/渠道容量约束。鲁棒优化中的不确定集与鲁棒预算 $\Gamma$ 取第二章口径；第三章滚动优化基线采用风险参数 $(\alpha,\lambda)=(0.95,0.6)$，切换惩罚系数 $\kappa$ 随年度分段递增以抑制跳变。对候选方案进行外层 $N=1000$ 次蒙特卡洛评估，统计关键指标并进行图示对比。

\subsection{评价指标}
本章采用以下核心评价指标：
\begin{itemize}
  \item \textbf{期望收益}：衡量整体回报水平；
  \item \textbf{最坏5\%分位}：反映尾部可保障程度；
  \item \textbf{$\mathrm{CVaR}_{95\%}$}：以 Rockafellar–Uryasev 形式衡量超越分位阈值的尾部期望损失；
  \item \textbf{后悔值}：定义为相对最优基准的收益让渡，用于比较不同保守度与风险规制下的效率损失；
  \item \textbf{资源利用指标}：峰值与平滑度反映对瓶颈资源的匹配能力与过载风险；
  \item \textbf{结构稳定性}：切换频率/幅度量化跨年度作物—地块配置的稳定性；
  \item \textbf{约束违约风险}：考察在不同保守度/风险参数下，容量、用工、水肥等约束的可行性边际。
\end{itemize}

\subsection{结果呈现与对比}
从收益分布看，确定性方案在波动条件下呈现较厚的尾部下行风险。相比之下，滚动 + CVaR 方案在让渡极少期望收益的同时，显著改善了最坏分位与 $\mathrm{CVaR}_{95\%}$ 指标，尾部收敛性能更优，如 \Cref{fig:eval-profit-dist} 所示。

\begin{figure}[H]
  \centering
  \begin{subfigure}[t]{0.48\linewidth}
    \centering
    \includegraphics[width=\linewidth]{model1_profit_distribution.pdf}
    \caption{确定性方案收益分布}
  \end{subfigure}\hfill
  \begin{subfigure}[t]{0.48\linewidth}
    \centering
    \includegraphics[width=\linewidth]{model3_profit_distribution_baseline.pdf}
    \caption{滚动框架对照基线}
  \end{subfigure}
  \caption{收益分布与尾部对比（确定性 vs 滚动基线）}
  \label{fig:eval-profit-dist}
\end{figure}

\FloatBarrier

鲁棒预算—收益前沿与后悔值对比显示：鲁棒优化与滚动 + CVaR 在稳健性目标上保持一致，但约束尾部风险的方式存在本质差异。具体而言：\begin{itemize}\item \textbf{鲁棒优化}：通过不确定集内最坏情景保证提供确定性保障，可解释性较强；\item \textbf{CVaR方法}：采用概率加权方式刻画尾部风险，更贴近现代风险管理实践。\end{itemize}两类方法在保守度—期望收益权衡曲线上形成相互可比的非支配解集，见 \Cref{fig:eval-frontier-regret}。

\begin{figure}[H]
  \centering
  \begin{subfigure}[t]{0.48\linewidth}
    \centering
    \includegraphics[width=\linewidth]{model2_gamma_frontier.pdf}
    \caption{鲁棒预算—收益前沿（第二章）}
  \end{subfigure}\hfill
  \begin{subfigure}[t]{0.48\linewidth}
    \centering
    \includegraphics[width=\linewidth]{model2_regret_curve.pdf}
    \caption{后悔值曲线（第二章）}
  \end{subfigure}
  \\[0.6em]
  \begin{subfigure}[t]{0.48\linewidth}
    \centering
    \includegraphics[width=\linewidth]{model3_gamma_frontier.pdf}
    \caption{鲁棒预算—收益前沿（第三章对照）}
  \end{subfigure}\hfill
  \begin{subfigure}[t]{0.48\linewidth}
    \centering
    \includegraphics[width=\linewidth]{model3_regret_curve.pdf}
    \caption{后悔值曲线（第三章对照）}
  \end{subfigure}
  \caption{鲁棒前沿与后悔值并排对比（第二章 vs 第三章）}
  \label{fig:eval-frontier-regret}
\end{figure}

\FloatBarrier

资源利用方面，鲁棒与滚动方案相较确定性方案能够有效削峰填谷，峰值更平缓、跨季更平滑，反映出对极端情景的更好适配能力，如 \Cref{fig:eval-util-compare} 所示。

\begin{figure}[H]
  \centering
  \begin{subfigure}[t]{0.32\linewidth}
    \centering
    \includegraphics[width=\linewidth]{model1_resource_util_line.pdf}
    \caption{确定性}
  \end{subfigure}
  \begin{subfigure}[t]{0.32\linewidth}
    \centering
    \includegraphics[width=\linewidth]{model2_resource_util_line_robust.pdf}
    \caption{鲁棒}
  \end{subfigure}
  \begin{subfigure}[t]{0.32\linewidth}
    \centering
    \includegraphics[width=\linewidth]{model3_resource_util_line_robust.pdf}
    \caption{滚动 + CVaR}
  \end{subfigure}
  \caption{资源利用率对比（确定性 vs 鲁棒 vs 滚动）}
  \label{fig:eval-util-compare}
\end{figure}

\FloatBarrier

结构演变与切换热图揭示：滚动 + CVaR 在容量瓶颈季节对高弹性作物进行分散配置，并在跨年间维持较低的切换幅度，兼顾收益与实施稳定性，见 \Cref{fig:eval-structure-switch}。

\begin{figure}[H]
  \centering
  \begin{subfigure}[t]{0.32\linewidth}
    \centering
    \includegraphics[width=\linewidth]{model1_stack_area.pdf}
    \caption{结构演变（确定性）}
  \end{subfigure}
  \begin{subfigure}[t]{0.32\linewidth}
    \centering
    \includegraphics[width=\linewidth]{model2_structure_stack_area_gamma.pdf}
    \caption{结构演变（鲁棒）}
  \end{subfigure}
  \begin{subfigure}[t]{0.32\linewidth}
    \centering
    \includegraphics[width=\linewidth]{model3_stack_area.pdf}
    \caption{结构演变（滚动）}
  \end{subfigure}
  \\[0.6em]
  \begin{subfigure}[t]{0.32\linewidth}
    \centering
    \includegraphics[width=\linewidth]{model1_switch_heatmap.pdf}
    \caption{切换热图（确定性）}
  \end{subfigure}
  \begin{subfigure}[t]{0.32\linewidth}
    \centering
    \includegraphics[width=\linewidth]{model2_switch_heatmap_robust.pdf}
    \caption{切换热图（鲁棒）}
  \end{subfigure}
  \begin{subfigure}[t]{0.32\linewidth}
    \centering
    \includegraphics[width=\linewidth]{model3_switch_heatmap.pdf}
    \caption{切换热图（滚动）}
  \end{subfigure}
  \caption{结构演变与切换稳定性对比}
  \label{fig:eval-structure-switch}
\end{figure}

\FloatBarrier

敏感性分析表明，与确定性和鲁棒方案相比，滚动 + CVaR 策略在以下方面表现更优：\begin{itemize}\item 显著降低了价格弹性参数对尾部指标的敏感度；\item 减少了容量参数变化对收益分布的影响；\item 在较高波动区间内提供了更强的约束可行性保障。\end{itemize}详细的违约风险对比分析见 \Cref{fig:eval-sensitivity-risk}。

\begin{figure}[H]
  \centering
  \begin{subfigure}[t]{0.32\linewidth}
    \centering
    \includegraphics[width=\linewidth]{model1_sensitivity_spider.pdf}
    \caption{敏感性（确定性）}
  \end{subfigure}
  \begin{subfigure}[t]{0.32\linewidth}
    \centering
    \includegraphics[width=\linewidth]{model2_sensitivity_spider_robust.pdf}
    \caption{敏感性（鲁棒）}
  \end{subfigure}
  \begin{subfigure}[t]{0.32\linewidth}
    \centering
    \includegraphics[width=\linewidth]{model3_sensitivity_spider_robust.pdf}
    \caption{敏感性（滚动）}
  \end{subfigure}
  \\[0.6em]
  \begin{subfigure}[t]{0.48\linewidth}
    \centering
    \includegraphics[width=\linewidth]{model2_violation_risk_bar.pdf}
    \caption{违约风险（鲁棒）}
  \end{subfigure}\hfill
  \begin{subfigure}[t]{0.48\linewidth}
    \centering
    \includegraphics[width=\linewidth]{model3_violation_risk_bar.pdf}
    \caption{违约风险（滚动）}
  \end{subfigure}
  \caption{敏感性与约束违约风险对比}
  \label{fig:eval-sensitivity-risk}
\end{figure}

\FloatBarrier

在价格—供给—容量反馈机制下，收益函数在容量约束附近表现出明显的凹性特征。滚动 + CVaR 策略能够：\begin{itemize}\item 有效识别并规避容量瓶颈处的边际收益骤降区间；\item 在高波动与销量受限的复杂场景下稳定提升保底收益；\item 通过动态调整实现风险控制与收益优化的平衡。\end{itemize}具体的收益—容量关系分析见 \Cref{fig:eval-revenue-capacity}。

\begin{figure}[H]
  \centering
  \begin{subfigure}[t]{0.32\linewidth}
    \centering
    \includegraphics[width=\linewidth]{model1_revenue_curve.pdf}
    \caption{收入—供给（确定性）}
  \end{subfigure}
  \begin{subfigure}[t]{0.32\linewidth}
    \centering
    \includegraphics[width=\linewidth]{model2_revenue_curve.pdf}
    \caption{收入—供给（鲁棒）}
  \end{subfigure}
  \begin{subfigure}[t]{0.32\linewidth}
    \centering
    \includegraphics[width=\linewidth]{model3_revenue_curve.pdf}
    \caption{收入—供给（滚动）}
  \end{subfigure}
  \\[0.6em]
  \begin{subfigure}[t]{0.32\linewidth}
    \centering
    \includegraphics[width=\linewidth]{model1_sales_vs_capacity.pdf}
    \caption{销量—产能（确定性）}
  \end{subfigure}
  \begin{subfigure}[t]{0.32\linewidth}
    \centering
    \includegraphics[width=\linewidth]{model3_sales_vs_capacity.pdf}
    \caption{销量—产能（滚动）}
  \end{subfigure}
  \begin{subfigure}[t]{0.32\linewidth}
    \centering
    \includegraphics[width=\linewidth]{model3_revenue_curve_capped.pdf}
    \caption{容量拐点与收入上限（滚动）}
  \end{subfigure}
  \caption{价格—供给—容量反馈与拐点效应}
  \label{fig:eval-revenue-capacity}
\end{figure}

\FloatBarrier

\subsection{讨论与管理启示}
综上，鲁棒优化与滚动 + CVaR 为管理者提供了互补的风险控制范式：前者以不确定集内最坏保证强化“确定性可解释性”，后者以概率加权的尾部约束贴近风险管理实践。价格弹性与容量拐点共同决定了在瓶颈区的边际收益递减，建议在年度计划中显式上报 CVaR@95\% 与最坏分位等尾部指标，并对地块切换设定年度“红线频率/幅度”，在保持实施稳定性的前提下实现收益—风险的透明调参（如 $\Gamma,\alpha,\lambda,\kappa$）。

\subsection{本章小结}
在统一口径的外层模拟评估下，滚动 + CVaR 相对确定性/鲁棒单策在相近期望收益下显著改善尾部指标，并降低对价格弹性与容量参数的尾部敏感度；资源曲线更平缓、结构更稳定。该框架为跨年决策提供了可操作的风险管理抓手，也为不同风险偏好的管理者提供了清晰的调参路径。
% ==================== 本章结束 ====================